% ══════════════════════════════════════════════════════════════
%  Código generado automáticamente por evaluate_classifier.py
%  Insertar en 8Resultados.tex dentro de \section{Evaluación...}
% ══════════════════════════════════════════════════════════════

\subsection{Evaluación sobre conjunto de prueba etiquetado}
\label{subsec:res_metricas_clasificador}

Para validar el desempeño del clasificador escalonado sobre datos \textbf{no vistos
durante el fine-tuning}, se construyó un conjunto de prueba de $n=41$ noticias
seleccionadas aleatoriamente del corpus de producción, excluyendo explícitamente
las instancias utilizadas en el entrenamiento del prototipo~4.
Las noticias fueron etiquetadas de forma independiente por dos anotadores
siguiendo un protocolo de codificación binaria:\
etiqueta $1$ para casos fácticos de orfandad por feminicidio y etiqueta $0$
para cualquier otro contenido.
De los $43$ artículos originales, $2$ fueron excluidos por falta de acuerdo entre anotadores.
El acuerdo inter-anotador se midió mediante el coeficiente $\kappa$ de
Cohen~\cite{cohen1960coefficient}, obteniéndose $\kappa = 0.91$, valor
que corresponde a un nivel de acuerdo \textit{casi perfecto} según la escala
de Landis y Koch~\cite{landis1977measurement}.

\vspace{0.5cm}
\begin{table}[htbp]
\centering
\caption{Métricas de evaluación del clasificador BETO escalonado sobre el conjunto de prueba
etiquetado ($n = 41$ noticias con consenso de anotadores)}
\label{tab:metricas_clasificador}
\renewcommand{\arraystretch}{1.3}
\begin{tabular}{|l|c|}
\hline
\rowcolor{black}
\textcolor{white}{\textbf{Métrica}} & \textcolor{white}{\textbf{Valor}} \\
\hline
Precisión (VP / VP+FP) & 0.77 \\
\hline
Recall --- sensibilidad (VP / VP+FN) & 0.95 \\
\hline
F1-Score & 0.85 \\
\hline
Tasa de Falsos Positivos & 30.0\% \\
\hline
$\kappa$ de Cohen (acuerdo inter-anotador) & 0.91 \\
\hline
\end{tabular}
\end{table}
\vspace{0.4cm}

\vspace{0.5cm}
\begin{table}[htbp]
\centering
\caption{Matriz de confusión del clasificador BETO escalonado (conjunto de prueba, $n=41$)}
\label{tab:confusion_matrix}
\renewcommand{\arraystretch}{1.4}
\begin{tabular}{|l|c|c|}
\hline
\rowcolor{black}
 & \textcolor{white}{\textbf{Predicción: No relevante}} &
   \textcolor{white}{\textbf{Predicción: Alta relevancia}} \\
\hline
\textbf{Gold: No relevante} & 14 (Verdadero Negativo) & 6 (Falso Positivo) \\
\hline
\textbf{Gold: Alta relevancia} & 1 (Falso Negativo) & 20 (Verdadero Positivo) \\
\hline
\end{tabular}
\end{table}
\vspace{0.4cm}
